
% \subsection{Summary of the architecture}
% \label{sec:method-summary}

% The proposed model is therefore characterized by two structural constraints:
% \begin{enumerate}
%     \item the encoder is exactly positively homogeneous of degree \(p\), ensuring explicit transport of regular variation to latent space;
%     \item the decoder is a homogeneous leading map plus a sparsity-regularized correction, allowing geometric flexibility at finite scales while retaining asymptotic homogeneity in the tail.
% \end{enumerate}
% This asymmetric treatment of encoder and decoder is central to the method: exact scaling control is needed to preserve latent tail structure, whereas only asymptotic scaling control is needed for reconstruction.



% % TODO: Write 2-3 paragraphs on motivation
% % Content to include:
% % - Heavy-tailed data is ubiquitous (finance, climate, networks, anomalies)
% % - Extreme value analysis needs to be done, but data is high-dimensional
% % - Dimensionality reduction is necessary, but standard methods don't preserve tail structure
% % - Example: financial crisis analysis requires understanding tail behavior in latent space
% \paragraph{Motivation.}
% TODO: Write compelling motivation for tail-preserving dimensionality reduction. Include:
% \begin{itemize}
%     \item Why heavy-tailed data matters (finance, climate, networks)
%     \item Need for dimensionality reduction in EVT
%     \item Failure of standard autoencoders to preserve tail structure
%     \item Concrete example or application
% \end{itemize}

% % TODO: Problem statement paragraph
% % Content: What are we trying to learn? Input/output, constraints, objectives
% \paragraph{Problem statement.}
% TODO: Clearly state the learning problem. Should include:
% \begin{itemize}
%     \item Input: Data X on manifold \(\cM \subset \RR^D\) with regularly varying distribution
%     \item Goal: Learn encoder \(f: \RR^D \to \RR^m\) (with \(m < D\)) that:
%     \begin{enumerate}
%         \item Preserves tail measure (for EVT in latent space)
%         \item Allows faithful reconstruction
%         \item Enables generation of extreme values
%     \end{enumerate}

% \end{itemize}

% % TODO: Why existing approaches fail
% \paragraph{Challenges.}
% TODO: Explain why this is hard and why standard methods fail. Include:
% \begin{itemize}
%     \item Standard autoencoders: no homogeneity → tail structure not preserved
%     \item Homogeneous maps alone: cannot handle non-homogeneous manifolds
%     \item Local vs asymptotic homogeneity: manifolds may be complex locally but cone-like in tail
%     \item Need for adaptive mechanism
% \end{itemize}

% % TODO: Our approach paragraph
% \paragraph{Our approach.}
% TODO: Describe our solution at high level (save details for Method section). Include:
% \begin{itemize}
%     \item Encoder: exactly \(p\)-homogeneous via product form \(w = r^p \cdot g(u)\)
%     \item Decoder: three components with adaptive homogeneity
%     \item Key innovation: sparsity-regularized correction term
%     \item Intuition: be flexible locally, homogeneous in tail
% \end{itemize}

% % TODO: Contributions list with quantitative claims where possible
% \paragraph{Contributions.}
% TODO: List contributions. For NeurIPS acceptance, should include:
% \begin{enumerate}[leftmargin=1.5em]
%     \item \textbf{Method}: Three-component decoder architecture with adaptive asymptotic homogeneity (Section~\ref{sec:method})
%     \item \textbf{Theory}: Analysis of tail preservation under bounded scaling and sparsity (Section~\ref{sec:theory})
%     \item \textbf{Experiments}: Validation on synthetic and real-world heavy-tailed data (Section~\ref{sec:experiments})
%     \begin{itemize}
%         \item TODO: Add quantitative claim (e.g., ``X\% improvement in tail index preservation'')
%         \item TODO: Add claim about generation quality
%     \end{itemize}
%     \item \textbf{Insights}: When does the sparsity mechanism work? Connection to manifold geometry (Section~\ref{sec:discussion})
%     \item \textbf{Code}: Open-source implementation TODO: add repo name/link
% \end{enumerate}

% % ==============================================================================
% % SECTION 2: BACKGROUND AND RELATED WORK
% % ==============================================================================
% % TODO LIST FOR THIS SECTION:
% % - [  ] Write Extreme Value Theory in ML paragraph (3-5 citations)
% % - [  ] Write Manifold Learning and Autoencoders paragraph (5-8 citations)
% % - [  ] Write Homogeneous Functions and Scaling paragraph (2-3 citations)
% % - [  ] Write Radial/Angular Decompositions paragraph (2-3 citations)
% % ==============================================================================

% \section{Related Work}
% \label{sec:related}

% % TODO: Write paragraph on EVT in ML
% \paragraph{Extreme Value Theory in Machine Learning.}
% The fields of Extreme Value Theory and Machine Learning are increasingly being considered together, with the goal of combining two superficially incompatible topics in data-driven analysis: machine learning methods rely on large datasets whilst EVT is primarily concerned with the data-sparse regions in the tails of distributions. Recent works such as \cite{Lafon_2023}, \cite{Allouche_2026}, \cite{Monte_2025}... use unsupervised learning techniques from the machine learning community alongside EVT principles to conduct generative modelling of multivariate extremes. 

% However little attention has been given to (nonlinear) dimension reduction. 

% Linear dimension reduction has been studied by...



% TODO: Discuss existing EVT applications in ML. Include:
% \begin{itemize}
%     \item Existing EVT methods (tail modeling, extreme quantiles)
%     \item Applications (anomaly detection, risk assessment)
%     \item Gap: no dimensionality reduction preserving tail structure
%     \item Citations: [TODO: add 3-5 papers on EVT in ML]
% \end{itemize}

% % TODO: Write paragraph on manifold learning
% \paragraph{Manifold Learning and Autoencoders.}
% TODO: Discuss representation learning for manifolds. Include:
% \begin{itemize}
%     \item Standard autoencoders, VAEs, \(\beta\)-VAEs
%     \item Geometric deep learning (hyperbolic, spherical embeddings)
%     \item Manifold learning methods (Isomap, LLE, etc.)
%     \item Limitation: none preserve heavy-tailed structure
%     \item Citations: [TODO: add 5-8 papers on autoencoders and manifold learning]
% \end{itemize}

% % TODO: Write paragraph on homogeneous functions
% \paragraph{Homogeneous Functions and Scaling.}
% TODO: Discuss role of homogeneity in EVT and representation learning. Include:
% \begin{itemize}
%     \item Regular variation and homogeneity (Resnick, de Haan \& Ferreira)
%     \item Scale-invariant representations
%     \item Our contribution: learning homogeneous maps for tail preservation
%     \item Citations: [TODO: add 2-3 EVT theory papers]
% \end{itemize}

% % TODO: Write paragraph on radial/angular decompositions
% \paragraph{Radial and Angular Decompositions.}
% TODO: Discuss related geometric decompositions. Include:
% \begin{itemize}
%     \item Radial basis functions, polar coordinates
%     \item Star-shaped parameterizations, radial graphs
%     \item Our contribution: adaptive asymptotic behavior via sparsity
%     \item Citations: [TODO: add 2-3 papers on radial parameterizations]
% \end{itemize}

% % ==============================================================================
% % SECTION 3: PROBLEM SETUP AND PRELIMINARIES
% % ==============================================================================
% % TODO LIST FOR THIS SECTION:
% % - [  ] Write data model paragraph (manifold M, regular variation)
% % - [  ] Write goal paragraph (encoder/decoder, requirements)
% % - [  ] Write key constraint paragraph (homogeneity for tail preservation)
% % - [  ] Add notation table if needed
% % ==============================================================================

% \section{Problem Setup and Preliminaries}
% \label{sec:setup}

% % TODO: Write data model paragraph
% \paragraph{Data Model.}

% TODO: Describe the data distribution and manifold structure. Include:
% \begin{itemize}
%     \item Data \(X \in \cM \subset \RR^D\) where \(\cM\) is a smooth \(m\)-dimensional manifold
%     \item Regular variation: \(t P(X/t \in A) \to \mu(A)\) as \(t \to \infty\)
%     \item Tail measure \(\mu\) encodes extreme behavior
%     \item Example: Pareto-distributed radius with uniform angular component
% \end{itemize}

% % TODO: Write goal paragraph
% \paragraph{Learning Goal.}
% TODO: State the learning problem precisely. Include:
% \begin{itemize}
%     \item Learn encoder \(f: \RR^D \to \RR^m\) and decoder \(g: \RR^m \to \RR^D\)
%     \item Requirements:
%     \begin{enumerate}[label=(\roman*)]
%         \item \textbf{Tail preservation}: latent tail measure \(\mu_W = f_* \mu_X\) (pushforward)
%         \item \textbf{Reconstruction}: \(\norm{X - g(f(X))}\) small on manifold
%         \item \textbf{Generation}: sample from tail in latent space, decode faithfully to extremes
%     \end{enumerate}
%     \item Challenge: simultaneous dimensionality reduction and tail preservation
% \end{itemize}

% % TODO: Write key constraint paragraph
% \paragraph{Key Constraint: Encoder Homogeneity.}
% TODO: Explain why homogeneity is necessary. Include:
% \begin{itemize}
%     \item Encoder must be \(p\)-homogeneous: \(f(\lambda x) = \lambda^p f(x)\)
%     \item \textbf{Theorem} (from EVT): If \(X\) is regularly varying and \(f\) is \(p\)-homogeneous, then \(W = f(X)\) is regularly varying with \(\mu_W = f_* \mu_X\)
%     \item Implication: can perform EVT analysis in latent space \(\RR^m\)
%     \item This constraint motivates our architectural design
% \end{itemize}

% % TODO: Add notation table if space allows
% % \paragraph{Notation.}
% % [Optional: table of key notation if space allows]

% % ==============================================================================
% % SECTION 4: METHOD - HOMOGENEOUS AUTOENCODERS
% % ==============================================================================
% % TODO LIST FOR THIS SECTION:
% % - [  ] Write encoder architecture subsection (decomposition, implementation, homogeneity proof)
% % - [  ] Write decoder challenge subsection (problem statement, why homogeneity is hard)
% % - [  ] Write three-component decoder subsection (architecture, forward pass, sparsity, intuition)
% % - [  ] Add architecture diagram (Figure 1 or 2)
% % - [  ] Add pseudocode for forward pass
% % ==============================================================================

% TODO: CHECK THIS SECTION
% \section{Method: Homogeneous Autoencoders}
% \label{sec:method}

% We now introduce the proposed autoencoder architecture. The design separates the roles of encoding and decoding. The encoder is constrained to be exactly \(p\)-homogeneous, so that latent regular variation follows from \Cref{prop:...}. The decoder is allowed to be non-homogeneous at moderate scales in order to be able to reconstruct curved manifold geometry, but is structured so that its non-homogeneous component is incentivised to be negligible in the tail, retaining asymptotic homogeneity and compatibility with transport of regular variation from latent to ambient spaces.

% \subsection{Encoder: \texorpdfstring{\(p\)}{p}-Homogeneous Representation}
% \label{sec:method:encoder}

% Fix a centre \(c\in\RR^D\) and norm indices \(q_r,q_\pi\geq 1\). For \(x\neq c\), define
% \[
% r_c(x):=\norm{x-c}_{q_r},
% \qquad
% \pi_c(x):=\frac{x-c}{\norm{x-c}_{q_\pi}}\in \mathbb{S}^{D-1}_{q_\pi},
% \]
% where
% \[
% \mathbb{S}^{D-1}_{q_\pi}:=\{u\in\RR^D:\norm{u}_{q_\pi}=1\}.
% \]
% Thus \(r_c\) is the radial variable used for scaling, while \(\pi_c\) is the angular projection used to parameterise direction.

% We define the encoder by
% \[
% f(x)=r_c(x)^p\, g(\pi_c(x)),
% \qquad x\neq c,
% \]
% where
% \[
% g:\mathbb{S}^{D-1}_{q_\pi}\to\RR^m
% \]
% is a measurable angular map. This form separates radial and angular information: the radial contribution enters only through \(r_c(x)^p\), whereas all directional variation is captured by \(g(\pi_c(x))\).

% A convenient further decomposition is
% \[
% g(u)=a(u)e(u),
% \]
% where
% \[
% a(u):=\norm{g(u)}_2\in(0,\infty),
% \qquad
% e(u):=\frac{g(u)}{\norm{g(u)}_2}\in\mathbb{S}^{m-1}.
% \]
% Hence, whenever \(g(u)\neq 0\),
% \[
% f(x)=r_c(x)^p\, a(\pi_c(x))\, e(\pi_c(x)).
% \]
% This introduces no additional restriction beyond nonvanishing of \(g\); it is simply the polar decomposition of the latent angular map \(g\) into a scalar magnitude term and a directional term.

% In implementation, we parameterize \(g\) either directly or via
% \[
% e(u)=\frac{E_{\mathrm{ang}}(u)}{\norm{E_{\mathrm{ang}}(u)}_2},
% \qquad
% a(u)=\mathrm{softplus}\!\bigl(E_{\mathrm{mag}}(u)\bigr),
% \]
% with neural networks
% \[
% E_{\mathrm{ang}}:\mathbb{S}^{D-1}_{q_\pi}\to\RR^m,
% \qquad
% E_{\mathrm{mag}}:\mathbb{S}^{D-1}_{q_\pi}\to\RR_+.
% \]

% The encoder is exactly \(p\)-homogeneous relative to the centre \(c\).

% \begin{lemma}
% \label{lem:encoder-homogeneous}
% Assume that the angular projection \(\pi_c\) is scale-invariant along rays through \(c\), i.e.
% \[
% \pi_c(c+\lambda(x-c))=\pi_c(x),
% \qquad x\neq c,\ \lambda>0.
% \]
% Then, for every \(x\neq c\) and every \(\lambda>0\),
% \[
% f\bigl(c+\lambda(x-c)\bigr)=\lambda^p f(x).
% \]
% In particular, when \(c=0\),
% \[
% f(\lambda x)=\lambda^p f(x),\qquad \lambda>0.
% \]
% \end{lemma}

% \begin{proof}
% For \(x\neq c\) and \(\lambda>0\),
% \[
% r_c(c+\lambda(x-c))=\lambda r_c(x)
% \]
% by homogeneity of the norm \(\norm{\cdot}_{q_r}\), and by assumption,
% \[
% \pi_c(c+\lambda(x-c))=\pi_c(x).
% \]
% Therefore
% \[
% f\bigl(c+\lambda(x-c)\bigr)
% =
% r_c(c+\lambda(x-c))^p\, g(\pi_c(c+\lambda(x-c)))
% =
% \lambda^p r_c(x)^p g(\pi_c(x))
% =
% \lambda^p f(x).
% \]
% \end{proof}

% By \Cref{prop:...}, this implies that if \(X\) is regularly varying and the remaining hypotheses there hold, then \(Z=f(X)\) is regularly varying with tail measure \(f_{\#}\mu\) and tail index \(\alpha/p\).

% \subsection{Decoder: Adaptive Asymptotic Homogeneity}
% \label{sec:method:decoder}

% Write the latent variable in polar form as
% \[
% z=\rho(z)\theta(z),
% \qquad
% \rho(z):=\norm{z}_2,
% \qquad
% \theta(z):=\frac{z}{\norm{z}_2}\in\mathbb{S}^{m-1}.
% \]
% Since the encoder scales ambient radius through \(r_c^p\), the decoder inverts this radial transformation through \(\rho^{1/p}\).

% We define the decoder using three components:
% \[
% c_{\mathrm{hom}}:\mathbb{S}^{m-1}\to\RR^D,
% \qquad
% b:\mathbb{S}^{m-1}\to(0,\infty),
% \qquad
% c_{\mathrm{cor}}:\mathbb{S}^{m-1}\times\RR\to\RR^D.
% \]
% Here \(c_{\mathrm{hom}}\) is the homogeneous angular component, \(b\) is an angle-dependent magnitude term, and \(c_{\mathrm{cor}}\) is a non-homogeneous correction. The function \(b\) depends only on latent angle and is the decoder analogue of the encoder-side magnitude modulation.

% The decoder is
% \[
% h(z)
% =
% c+\rho(z)^{1/p}
% \Bigl(
% b(\theta(z))\,c_{\mathrm{hom}}(\theta(z))
% +
% c_{\mathrm{cor}}(\theta(z), \rho(z))
% \Bigr).
% \]
% The correction is added at the end. Thus the homogeneous structure is explicit, while the non-homogeneous component enters as an additive perturbation.

% In implementation, we use
% \[
% c_{\mathrm{hom}}(\theta)
% =
% \frac{H_{\mathrm{hom}}(\theta)}{\norm{H_{\mathrm{hom}}(\theta)}_2},
% \qquad
% b(\theta)=\mathrm{softplus}\!\bigl(H_{\mathrm{mag}}(\theta)\bigr),
% \qquad
% c_{\mathrm{cor}}(\theta,\rho)=H_{\mathrm{cor}}(\theta,\rho),
% \]
% with
% \[
% H_{\mathrm{hom}}:\mathbb{S}^{m-1}\to\RR^D,
% \qquad
% H_{\mathrm{mag}}:\mathbb{S}^{m-1}\to\RR,
% \qquad
% H_{\mathrm{cor}}:\mathbb{S}^{m-1}\times\RR\to\RR^D.
% \]

% The decoder is trained jointly with the encoder using reconstruction loss plus a radius-weighted sparsity penalty on the correction term. For \(z=f(x)\), we use
% \[
% \mathcal{L}(x)
% =
% \mathcal{L}_{\mathrm{recon}}\!\bigl(x,h(f(x))\bigr)
% +
% \lambda_{\mathrm{cor}}\,
% \rho(f(x))\,
% \norm{c_{\mathrm{cor}}(\theta(f(x)), \rho(f(x)))}_1.
% \]
% The factor \(\rho(f(x))\) increases the cost of non-homogeneous corrections at large latent radii. Consequently, the decoder is encouraged to use the correction term at moderate scales, where it improves reconstruction, and suppress it in the tail, where asymptotic homogeneity is preferred.

% This is the basic mechanism of adaptive asymptotic homogeneity: local geometric flexibility is carried by \(c_{\mathrm{cor}}\), while tail behaviour is controlled by the homogeneous component \(b(\theta)c_{\mathrm{hom}}(\theta)\).
% % TODO: Write encoder subsection
% \subsection{Encoder: \texorpdfstring{\(p\)}{p}-Homogeneity}
% \label{sec:method:encoder}

% % TODO: Write decomposition paragraph
% \paragraph{Product-Form Decomposition.}
% TODO: Explain encoder construction. Include:
% \begin{itemize}
%     \item Decompose \(x = c + r \cdot u\) where \(r = \norm{x - c}_2\) and \(u = (x - c)/r \in \mathbb{S}^{D-1}\)
%     \item Product form: \(w = r^p \cdot a(u) \cdot e(u)\) where:
%     \begin{itemize}
%         \item \(e(u): \mathbb{S}^{D-1} \to \mathbb{S}^{m-1}\) (angular encoding, normalized to sphere)
%         \item \(a(u): \mathbb{S}^{D-1} \to \RR^+\) (magnitude scaling, bounded on compact domain)
%     \end{itemize}
%     \item Latent \(w \in \RR^m\) with polar structure: \(r_w = \norm{w}\), \(\theta = w/\norm{w} \in \mathbb{S}^{m-1}\)
% \end{itemize}

% % TODO: Write implementation paragraph
% \paragraph{Implementation.}
% TODO: Describe neural architecture. Include:
% \begin{itemize}
%     \item Two MLPs: \texttt{encoder\_angular} (\(\mathbb{S}^{D-1} \to \RR^m\), then normalize) and \texttt{encoder\_magnitude} (\(\mathbb{S}^{D-1} \to \RR^+\), apply softplus)
%     \item Forward pass: \(e = \text{normalize}(\text{MLP}_e(u))\), \(a = \text{softplus}(\text{MLP}_a(u))\), \(w = r^p \cdot a \cdot e\)
%     \item Center \(c\) is fixed (could be learned, currently origin or data center)
% \end{itemize}

% % TODO: Write homogeneity proof
% \paragraph{Homogeneity and Tail Preservation.}
% TODO: Prove encoder is \(p\)-homogeneous. Include:
% \begin{itemize}
%     \item For \(\lambda > 0\): \(u(\lambda x) = u(x)\) (direction unchanged), \(r(\lambda x) = \lambda r(x)\)
%     \item Therefore: \(\text{encode}(\lambda x) = (\lambda r)^p \cdot a(u) \cdot e(u) = \lambda^p \cdot \text{encode}(x)\)
%     \item \textbf{Consequence}: Tail measure preserved via pushforward: \(\mu_W = \text{encode}_* \mu_X\)
%     \item This enables EVT analysis in latent space
% \end{itemize}

% % TODO: Write decoder challenge subsection
% \subsection{Decoder Challenge: Topological Collapse}
% \label{sec:method:decoder-challenge}

% % TODO: Write problem statement
% \paragraph{The Problem.}
% TODO: Explain decoder difficulty. Include:
% \begin{itemize}
%     \item Latent \(w \in \RR^m\) has \(m\) degrees of freedom: radius \(r_w = \norm{w}\) and angle \(\theta = w/\norm{w} \in \mathbb{S}^{m-1}\)
%     \item Manifold \(\cM\) has dimension \(m\) (same!)
%     \item Na\"ive decoder \(c(\theta): \mathbb{S}^{m-1} \to \mathbb{S}^{D-1}\) uses only \((m-1)\) DoF (angle)
%     \item Result: decoder output collapses to \((m-1)\)-dimensional curve, cannot span \(m\)-dimensional manifold
% \end{itemize}

% % TODO: Write why homogeneity is hard
% \paragraph{Why Homogeneity is Hard for Decoders.}
% TODO: Explain the fundamental tension. Include:
% \begin{itemize}
%     \item Ideally: \(\text{decode}(\lambda^p w) = \lambda \cdot \text{decode}(w)\) for faithful generation from tail
%     \item But: manifolds are not necessarily homogeneous everywhere (complex local structure)
%     \item Encoder can be locally homogeneous (radial projection well-defined off-manifold)
%     \item Decoder cannot: scaling in latent space must map to manifold, which may not be cone-like locally
%     \item Solution: \emph{adaptive asymptotic homogeneity} via sparsity regularization
% \end{itemize}

% % TODO: Write three-component decoder subsection
% \subsection{Three-Component Decoder with Adaptive Homogeneity}
% \label{sec:method:decoder}

% % TODO: Write architecture overview
% \paragraph{Architecture Overview.}
% TODO: Describe three-component structure. Include:
% \begin{enumerate}
%     \item \(c_{\text{base}}(\theta): \mathbb{S}^{m-1} \to \mathbb{S}^{D-1}\) --- homogeneous baseline angular component
%     \item \(c_{\text{correction}}(\theta, r_w): \mathbb{S}^{m-1} \times \RR \to \RR^D\) --- sparse non-homogeneous correction
%     \item \(b(\theta): \mathbb{S}^{m-1} \to \RR^+\) --- bounded angle-dependent radius scaling
% \end{enumerate}
% Key insight: separate homogeneous baseline from non-homogeneous correction, penalize correction to drive asymptotic homogeneity.

% % TODO: Write forward pass
% \paragraph{Forward Pass.}
% TODO: Describe decoding procedure. Include pseudocode:
% \begin{verbatim}
% r_w = ||w||
% theta = w / ||w||  (angle in latent space)
% r_reconstructed = r_w^(1/p)  (invert encoder's r^p)

% c_base = normalize(MLP_base(theta))  (homogeneous baseline)
% c_correction = MLP_correction([theta, (r_w)])  (non-homogeneous)
% c = normalize(c_base + c_correction)  (combined direction)

% b = softplus(MLP_magnitude(theta))  (bounded scaling)
% h = b * c  (final radial component)
% x_hat = center + r_reconstructed * h  (ambient reconstruction)
% \end{verbatim}

% % TODO: Write sparsity regularization
% \paragraph{Sparsity Regularization.}
% TODO: Explain loss function and mechanism. Include:
% \begin{itemize}
%     \item Total loss: \(\cL = \cL_{\text{recon}} + \lambda_{\text{ang}} \cL_{\text{ang}} + \lambda_{\text{rad}} \cL_{\text{rad}} + \lambda_{\text{sparse}} \norm{c_{\text{correction}}}_1?*r_w?\)
%     \item L1 penalty on \(c_{\text{correction}}\) and weighted by r encourages sparsity (entries driven to zero when not needed)
%     \item Trade-off: reconstruction quality vs asymptotic homogeneity
%     \item Mechanism: when \(c_{\text{base}}\) alone suffices, correction is penalized to zero
% \end{itemize}

% % TODO: Write intuition paragraph
% \paragraph{Intuition: Adaptive Homogeneity.}
% TODO: Explain how this solves the tension. Include:
% \begin{itemize}
%     \item \textbf{Near origin} (small \(r_w\)): manifold may have complex, non-homogeneous structure
%     \begin{itemize}
%         \item Reconstruction loss dominates sparsity penalty
%         \item \(c_{\text{correction}}\) can be non-zero to capture complexity
%         \item Decoder is locally flexible
%     \end{itemize}
%     \item \textbf{In tail} (large \(r_w\)): manifold becomes asymptotically cone-like (due to regular variation)
%     \begin{itemize}
%         \item \(c_{\text{base}}(\theta)\) alone captures structure
%         \item \(c_{\text{correction}}\) provides diminishing benefit
%         \item Sparsity penalty drives \(\norm{c_{\text{correction}}} \to 0\)
%         \item Decoder becomes effectively homogeneous: \(c \approx c_{\text{base}}(\theta)\)
%     \end{itemize}
%     \item Result: automatic adaptation to manifold geometry ADD NOTES ABOUT WEIGHTING BY RADIUS
% \end{itemize}

% % TODO: Add architecture diagram reference
% % [TODO: Add Figure 2 - Architecture diagram showing encoder/decoder components]

% % ==============================================================================
% % SECTION 5: THEORETICAL ANALYSIS
% % ==============================================================================
% % TODO LIST FOR THIS SECTION:
% % - [  ] Write encoder tail preservation theorem + proof sketch
% % - [  ] Write decoder tail index preservation proposition + analysis
% % - [  ] Write sparsity induces asymptotic homogeneity proposition + trade-off analysis
% % - [  ] Add brief connection to geometric theory (defer full theory to appendix)
% % ==============================================================================

% \section{Theoretical Analysis}
% \label{sec:theory}

% % TODO: State theorem
% \begin{theorem}[Encoder preserves tail measure]
% \label{thm:encoder-tail-preservation}
% Let \(X \in \RR^D\) be regularly varying with tail measure \(\mu_X\). If the encoder \(f: \RR^D \to \RR^m\) is \(p\)-homogeneous (i.e., \(f(\lambda x) = \lambda^p f(x)\) for all \(\lambda > 0\)), then \(W = f(X)\) is regularly varying with tail measure \(\mu_W = f_* \mu_X\) (the pushforward of \(\mu_X\)).
% \end{theorem}

% % TODO: Proof sketch
% \begin{proof}[Proof sketch]
% TODO: Provide proof sketch. Include:
% \begin{itemize}
%     \item Regular variation: \(t P(X/t \in A) \to \mu_X(A)\) as \(t \to \infty\)
%     \item Apply \(f\): \(t P(f(X)/t^p \in B) = t P(X/t \in f^{-1}(B)) \to \mu_X(f^{-1}(B))\)
%     \item Change of variables: \(s = t^p\), so \(s P(f(X)/s \in B) \to \mu_X(f^{-1}(B))\)
%     \item Define \(\mu_W(B) := \mu_X(f^{-1}(B))\), which is exactly the pushforward \(f_* \mu_X\)
%     \item Conclusion: EVT can be performed in latent space with tail measure exactly preserved
% \end{itemize}
% \end{proof}

% % TODO: Write decoder tail index preservation
% \subsection{Decoder: Tail Index Preservation via Bounded Scaling}
% \label{sec:theory:decoder}

% % TODO: State proposition
% \begin{proposition}[Bounded scaling preserves tail index]
% \label{prop:bounded-scaling}
% Suppose the decoder magnitude component \(b: \mathbb{S}^{m-1} \to \RR^+\) is bounded: \(b_{\min} \le b(\theta) \le b_{\max}\) for all \(\theta \in \mathbb{S}^{m-1}\). Then for generated samples \(W_{\text{new}}\) from the tail distribution in latent space, the decoded samples \(\hat{X} = \text{decode}(W_{\text{new}})\) preserve the tail index.
% \end{proposition}

% % TODO: Analysis
% \begin{proof}[Analysis]
% TODO: Provide analysis. Include:
% \begin{itemize}
%     \item Decoder output: \(\hat{X} = c + r_w^{1/p} \cdot b(\theta) \cdot \hat{c}\) where \(\hat{c} \in \mathbb{S}^{D-1}\)
%     \item Radial component: \(\norm{\hat{X} - c} = r_w^{1/p} \cdot b(\theta) \cdot \norm{\hat{c}} = r_w^{1/p} \cdot b(\theta)\)
%     \item Bounded \(b\): \(b_{\min} \cdot r_w^{1/p} \le \norm{\hat{X} - c} \le b_{\max} \cdot r_w^{1/p}\)
%     \item Asymptotic behavior: \(\norm{\hat{X} - c} \sim r_w^{1/p}\) as \(r_w \to \infty\)
%     \item Tail index: if \(P(\norm{W} > s) \sim s^{-\alpha}\), then \(P(\norm{\hat{X}} > t) \sim P(r_w^{1/p} > t/b_{\max}) \sim (t/b_{\max})^{-\alpha p} = t^{-\alpha p}\)
%     \item Conclusion: tail index transformed by \(p\) (as expected from \(p\)-homogeneity)
% \end{itemize}
% \end{proof}

% % TODO: Write sparsity mechanism
% \subsection{Sparsity Induces Asymptotic Homogeneity}
% \label{sec:theory:sparsity}

% % TODO: State proposition
% \begin{proposition}[Sparsity drives asymptotic homogeneity]
% \label{prop:sparsity-asymptotic}
% Under the L1 sparsity penalty \(\lambda_{\text{sparse}} \norm{c_{\text{correction}}}_1\), the correction term \(c_{\text{correction}}\) is driven to zero in regions where the homogeneous baseline \(c_{\text{base}}\) alone provides sufficient reconstruction quality.
% \end{proposition}

% % TODO: Intuitive argument
% \begin{proof}[Intuitive argument]
% TODO: Provide intuition. Include:
% \begin{itemize}
%     \item Loss trade-off: \(\cL = \cL_{\text{recon}} + \lambda_{\text{sparse}} \norm{c_{\text{correction}}}_1\)
%     \item When \(c_{\text{base}}\) sufficient: increasing \(\norm{c_{\text{correction}}}\) provides no reconstruction benefit but increases penalty
%     \item Gradient descent drives \(c_{\text{correction}} \to 0\)
%     \item In tail: manifold asymptotically cone-like (due to regular variation) \(\Rightarrow\) \(c_{\text{base}}\) sufficient
%     \item Result: \(c \approx c_{\text{base}}(\theta)\) for large \(r_w\), decoder becomes asymptotically homogeneous
%     \item Angular behavior: \(\hat{X}/\norm{\hat{X}} \approx c_{\text{base}}(\theta)\) depends only on latent angle, not radius
% \end{itemize}
% \end{proof}

% % TODO: Trade-off analysis
% \paragraph{Trade-off Analysis.}
% TODO: Discuss hyperparameter \(\lambda_{\text{sparse}}\). Include:
% \begin{itemize}
%     \item Small \(\lambda_{\text{sparse}}\): correction active everywhere, flexible but may not be asymptotically homogeneous
%     \item Large \(\lambda_{\text{sparse}}\): correction suppressed, asymptotically correct but may lose local flexibility
%     \item Optimal \(\lambda_{\text{sparse}}\): depends on manifold geometry (how quickly it becomes cone-like)
%     \item Manifold-dependent: complex manifolds need smaller \(\lambda\), simple cones need larger \(\lambda\)
% \end{itemize}

% % TODO: Connection to geometric theory
% \subsection{Connection to Geometric Theory}
% \label{sec:theory:geometry}

% % TODO: Brief mention
% \paragraph{Ray-Fiber Geometry (Brief).}
% TODO: Briefly mention geometric theory, defer to appendix. Include:
% \begin{itemize}
%     \item Full geometric theory (ray-fibers, immersion, injectivity) developed in Appendix~\ref{app:geometry}
%     \item Key insight: decoder must span full manifold dimension (immersion) and avoid collisions (injectivity)
%     \item Three-component decoder addresses both: correction term provides degrees of freedom when needed
%     \item Sparsity ensures correction not overused, maintaining interpretability and generalization
%     \item For interested readers: see Appendix for complete differential-geometric analysis
% \end{itemize}

% % ==============================================================================
% % SECTION 6: EXPERIMENTS
% % ==============================================================================
% % TODO LIST FOR THIS SECTION:
% % - [  ] Run and report Experiment 1: Perfect cone (homogeneous manifold)
% % - [  ] Run and report Experiment 2: Non-homogeneous surface
% % - [  ] Run and report Experiment 3: Ablation study (4 variants)
% % - [  ] Create visualization figures (sparsity, latent space, generation)
% % - [  ] Select 3 real-world datasets with heavy tails
% % - [  ] Implement 3 baseline methods (standard AE, VAE, β-VAE)
% % - [  ] Report tail preservation metrics, quantile prediction, generation quality
% % - [  ] Create Tables 1-3 and Figures 3-6
% % ==============================================================================

% \section{Experiments}
% \label{sec:experiments}

% % TODO: Write synthetic experiments subsection
% \subsection{Controlled Synthetic Experiments}
% \label{sec:experiments:synthetic}

% % TODO: Experiment 1
% \paragraph{Experiment 1: Perfect Cone (Homogeneous Manifold).}
% TODO: Describe cone experiment and results. Include:
% \begin{itemize}
%     \item \textbf{Setup}: Generate cone \(\{(r \cos \phi, r \sin \phi, hr) : r \sim \text{Pareto}(\alpha), \phi \sim \text{Uniform}(0, 2\pi)\}\)
%     \item Train with \(\lambda_{\text{sparse}} \in \{0.001, 0.01, 0.1\}\)
%     \item \textbf{Metrics}: sparsity \(\norm{c_{\text{correction}}}_1\), reconstruction MSE, tail index error
%     \item \textbf{Expected}: Low sparsity (correction not needed), excellent tail preservation
%     \item \textbf{Results}: [TODO: fill in after running experiments]
%     \item \textbf{Figure 3(a)}: Sparsity evolution during training
% \end{itemize}

% % TODO: Experiment 2
% \paragraph{Experiment 2: Non-Homogeneous Surface.}
% TODO: Describe surface experiment and results. Include:
% \begin{itemize}
%     \item \textbf{Setup}: Complex surface (current implementation)
%     \item Train with same \(\lambda_{\text{sparse}}\) values
%     \item \textbf{Expected}: Higher sparsity (correction active), still preserves tail index
%     \item \textbf{Results}: [TODO: fill in]
%     \item \textbf{Figure 3(b)}: Compare sparsity vs reconstruction trade-off
% \end{itemize}

% % TODO: Experiment 3
% \paragraph{Experiment 3: Ablation Study.}
% TODO: Report ablation results. Include:
% \begin{itemize}
%     \item \textbf{Variants tested}:
%     \begin{enumerate}
%         \item Full model (baseline + correction + sparsity)
%         \item No correction (baseline only, \(c_{\text{correction}} = 0\))
%         \item No sparsity (\(\lambda_{\text{sparse}} = 0\), correction always active)
%         \item Standard autoencoder (no homogeneity constraint)
%     \end{enumerate}
%     \item \textbf{Results}: [TODO: Table 1 - quantitative comparison]
%     \item \textbf{Figure 4}: Bar chart comparing reconstruction, tail index preservation, generation quality
% \end{itemize}

% % TODO: Visualization subsection
% \subsection{Visualization and Analysis}
% \label{sec:experiments:visualization}

% % TODO: List figures to create
% TODO: Create visualization figures:
% \begin{itemize}
%     \item \textbf{Figure 5}: \(\norm{c_{\text{correction}}}\) vs \(r_w\) (should decrease for large \(r_w\))
%     \item \textbf{Figure 6}: Latent space structure (2D scatter, colored by radius)
%     \item \textbf{Figure 7}: Generated extremes from tail (sample large \(r_w\), decode, compare with true extremes)
% \end{itemize}

% % TODO: Real-world experiments subsection
% \subsection{Real-World Heavy-Tailed Data}
% \label{sec:experiments:realworld}

% % TODO: Dataset selection
% \paragraph{Datasets.}
% TODO: Select and describe 3 datasets with documented heavy tails:
% \begin{enumerate}
%     \item \textbf{Financial returns}: [TODO: specify dataset, e.g., S\&P 500 daily returns]
%     \item \textbf{Network traffic}: [TODO: specify dataset]
%     \item \textbf{Climate/Insurance/Earthquake}: [TODO: select one and specify]
% \end{enumerate}
% For each: describe preprocessing, verify heavy tails (Hill estimator + QQ plots), train/val/test split.

% % TODO: Baselines
% \paragraph{Baselines.}
% TODO: Implement and compare with:
% \begin{enumerate}
%     \item Standard autoencoder (no homogeneity)
%     \item VAE (variational autoencoder)
%     \item \(\beta\)-VAE (with \(\beta > 1\) for better disentanglement)
%     \item Our method (with tuned \(\lambda_{\text{sparse}}\))
% \end{enumerate}

% % TODO: Metrics and results
% \paragraph{Metrics and Results.}
% TODO: Report results. Include:
% \begin{itemize}
%     \item \textbf{Table 2}: Tail index preservation (\(\alpha_{\text{original}}\), \(\alpha_{\text{latent}}\), \(\alpha_{\text{recon}}\), \(|\Delta \alpha|\))
%     \item \textbf{Table 3}: Extreme quantile prediction error (99\%, 99.9\%, 99.99\% quantiles)
%     \item \textbf{Figure 8}: QQ plots showing tail preservation (our method vs baselines)
%     \item \textbf{Figure 9}: Generated extremes vs real extremes (scatter plots)
%     \item \textbf{Key result}: Our method achieves [TODO: X\%] better tail index preservation than baselines while maintaining competitive reconstruction MSE
% \end{itemize}

% % ==============================================================================
% % SECTION 7: DISCUSSION
% % ==============================================================================
% % TODO LIST FOR THIS SECTION:
% % - [  ] Write when does it work paragraph
% % - [  ] Write hyperparameter selection paragraph
% % - [  ] Write limitations paragraph
% % - [  ] Write broader impact paragraph
% % ==============================================================================

% \section{Discussion}
% \label{sec:discussion}

% % TODO: When does it work
% \paragraph{When Does the Method Work?}
% TODO: Discuss conditions for success. Include:
% \begin{itemize}
%     \item Works well when: manifold is asymptotically cone-like (regular variation holds)
%     \item Struggles when: manifold has persistent curvature at all scales
%     \item Manifold complexity determines optimal \(\lambda_{\text{sparse}}\)
%     \item Empirical observation: sparsity metric correlates with manifold geometry
% \end{itemize}

% % TODO: Hyperparameter selection
% \paragraph{Hyperparameter Selection.}
% TODO: Provide guidance. Include:
% \begin{itemize}
%     \item \(\lambda_{\text{sparse}}\): tune based on manifold complexity (start with 0.01)
%     \item Monitor sparsity during training (should stabilize)
%     \item Could be adaptive in future work (increase with \(r_w\))
%     \item Other hyperparameters (\(p\), network architecture): see Appendix~\ref{app:implementation}
% \end{itemize}

% % TODO: Limitations
% \paragraph{Limitations.}
% TODO: Discuss limitations honestly. Include:
% \begin{itemize}
%     \item Requires choosing center \(c\) (currently fixed, could be learned)
%     \item Assumes manifold structure (dimension \(m\)) is known or estimated
%     \item Current implementation: 2D latent space (extension to higher \(m\) is future work)
%     \item Sparsity mechanism is implicit (could be more explicitly controlled)
% \end{itemize}

% % TODO: Broader impact
% \paragraph{Broader Impact.}
% TODO: Discuss societal impact. Include:
% \begin{itemize}
%     \item \textbf{Positive}: enables EVT in latent space, applications in risk assessment, climate modeling, anomaly detection
%     \item \textbf{Potential misuse}: could generate adversarial extremes (acknowledge this risk)
%     \item \textbf{Recommendations}: use responsibly, validate generated samples, consider ethical implications
% \end{itemize}

% % ==============================================================================
% % SECTION 8: CONCLUSION
% % ==============================================================================
% % TODO LIST FOR THIS SECTION:
% % - [  ] Write 2-3 sentence summary
% % - [  ] Write 1-2 sentence impact statement
% % - [  ] List 2-3 future work directions
% % ==============================================================================

% \section{Conclusion}
% \label{sec:conclusion}

% % TODO: Summary
% TODO: Summarize paper (2-3 sentences). Include:
% \begin{itemize}
%     \item Introduced homogeneous autoencoders with adaptive asymptotic homogeneity
%     \item Provably preserve tail structure through encoder homogeneity and sparsity-regularized decoder
%     \item Validated on synthetic and real-world heavy-tailed data with [TODO: X\%] improvement over baselines
% \end{itemize}

% % TODO: Impact
% TODO: State impact (1-2 sentences). Include:
% \begin{itemize}
%     \item Enables extreme value analysis in latent space
%     \item Opens new possibilities for high-dimensional EVT and generation of extreme events
% \end{itemize}

% % TODO: Future work
% TODO: Future directions:
% \begin{itemize}
%     \item Adaptive \(\lambda_{\text{sparse}}\) based on local manifold geometry and \(r_w\)
%     \item Extension to higher-dimensional latent spaces (\(m \> 2\))
%     \item Learning center \(c\) jointly with encoder/decoder
%     \item Applications to specific domains (finance, climate, networks)
%     \item Integration with normalizing flows for more flexible generation
% \end{itemize}

% \section*{References}

{\small
\bibliographystyle{plainnat}
\bibliography{references}
}

% ##############################################################################
% ##  APPENDIX SKELETON (uncomment \appendix and the sections below before    ##
% ##  submission; see MASTER PLAN §A4 for the six-appendix layout)            ##
% ##############################################################################
% TODO APX-0: Uncomment \appendix AND uncomment \section{...} headings for
% each appendix A-F. Each appendix has its own TODO block below.
%
% \appendix
%
% ==============================================================================
% APPENDIX A -- PROOFS
% ==============================================================================
% TODO APX-A1: Proof of Proposition 1 (encoder tail transport). Use the
% current proof at lines ~595-665 verbatim; cosmetic clean-up only.
% TODO APX-A2: Proof of Theorem 2 (decoder tail equivalence). Sketch
% provided in M-THM2a/b; write out fully. Length budget: 1 page.
% TODO APX-A3: Proof of Lemma 3.X (uniform angular convergence), if
% stated in §3.3. Short (half page).
%
% ==============================================================================
% APPENDIX B -- CANONICAL RADIAL-ANGULAR FACTORISATION (FULL)
% ==============================================================================
% TODO APX-B1: Move the entire \subsection{General form of homogeneous maps}
% here (see I-B1, I-B2). Keep Prop canonical-homogeneous-form + full proof
% + the a/e secondary decomposition + the general-radial-function remark.
% TODO APX-B2: Add the "Exact vs bounded radial preservation" Remark
% developed in I-B4 (three tiers: a=1, bounded a, unrestricted a).
%
% ==============================================================================
% APPENDIX C -- GEOMETRIC THEORY: RAY-FIBERS, IMMERSION, INJECTIVITY
% ==============================================================================
% TODO APX-C0: The user explicitly wants ray-fiber/immersion/injectivity
% theory in the appendix. Source material lives in the old commented-out
% \subsection{Ray-Fiber Analysis}, \subsection{Immersion Criteria},
% \subsection{Injectivity and Global Collisions} blocks further down (use
% as seed content). Consolidate into a coherent half-appendix (~2-3 pp).
% TODO APX-C1: \subsection{Ray-fibers.} Define F_c(u), S_c(u), Q_c(u,v).
% State the hierarchy (singleton, finite, discrete, positive-dim).
% TODO APX-C2: \subsection{Immersion criteria.} State and prove when
% d(f_c)_x is injective via the angular kernel K_x analysis.
% TODO APX-C3: \subsection{Injectivity.} Collision criterion and role of
% the magnitude term.
%
% ==============================================================================
% APPENDIX D -- IMPLEMENTATION AND HYPERPARAMETERS
% ==============================================================================
% TODO APX-D1: \subsection{Network architectures.} Exact layer sizes for
% A_phi, E_phi, A_psi, E_psi, Delta_psi. Source: experiments/lib/models.py
% (HomogeneousAutoencoder, StandardAutoencoder, PCABaseline) and the
% compute_matched_hidden_dim parameter-matching routine. Report total
% parameter counts for each model per experiment.
% TODO APX-D2: \subsection{Training.} Optimiser (Adam), LR, schedule,
% batch size, epochs, seeds, early-stopping. Source: experiments/lib/config.py
% (TrainConfig defaults) and experiments/climate/run_all.sh for climate.
% TODO APX-D3: \subsection{Adaptive lambda_cor schedule.} Full two-phase
% schedule (warmup T_A, EMA tolerance tau, growth g, decay d, cap lambda_max).
% Pseudocode preferred. Source: experiments/lib/evaluation.py if implemented
% there, else the training loop in lib.
% TODO APX-D4: \subsection{Metrics.} Exact definitions of hill_estimate
% (Hill 1975), hill_drift, extreme_quantile_errors, angular_tail_distance,
% encoder_homogeneity_error, tail_conditional_mse, extrapolation_mse,
% binned_reconstruction_error. Source: experiments/lib/metrics.py + the
% README at experiments/lib/README.md.
% TODO APX-D5: \subsection{Compute and reproducibility.} Hardware used,
% wall-clock per experiment, total GPU-hours, RNG seeds. Required for
% checklist Q8.
%
% ==============================================================================
% APPENDIX E -- ADDITIONAL EXPERIMENTS
% ==============================================================================
% TODO APX-E1: Homogeneity residual scan -- exp01/fig_homogeneity_scan.png.
% TODO APX-E2: HAE correction diagnostics -- exp02/fig_hae_correction_*.png.
% TODO APX-E3: Full per-sweep panels -- exp03 (D), exp04 (m), exp05 (alpha),
% exp06 (p), exp07 (hidden_dim). Use the fuller per-metric layouts.
% TODO APX-E4: Full extrapolation curves -- exp02/fig_extrapolation.png.
% TODO APX-E5: Climate PCA scree curves -- climate/results/<var>/pca_scree.png
% (four variants: u10, tp, t2m, joint).
% TODO APX-E6: Climate per-model spatial plots + marginal distributions
% (the subset not shown in main §5.3).
%
% ==============================================================================
% APPENDIX F -- EXTREME VALUE THEORY BACKGROUND
% ==============================================================================
% TODO APX-F1: \subsection{Regular variation.} Full definition,
% b-function, vague convergence in M(O_D), alpha-homogeneity. Cite
% Resnick 2007, de Haan & Ferreira 2006, Lindskog-Resnick-Roy 2014.
% Target audience: ML reviewer with no EVT background. ~1 page.
% TODO APX-F2: \subsection{Mapping theorem (full statement).} Formal
% statement as Theorem F.1, with proof pointer to lindskog_regularly_2014.
% TODO APX-F3: \subsection{Hill estimator.} Definition, choice of k, Hall
% & Welsh 1985 for bias-variance trade-off, our default (k = 10% of n).
% Cite Hill 1975 (ADD to bib).
% TODO APX-F4: Brief demonstration per I-M6a: exhibit a smooth bounded
% encoder f : R^D -> R^m that annihilates the tail measure, to make the
% "arbitrary nonlinear maps destroy regular variation" point rigorous in
% one paragraph.
% ##############################################################################
% ##                         END APPENDIX SKELETON                            ##
% ##############################################################################

% %%%%%%%%%%%%%%%%%%%%%%%%%%%%%%%%%%%%%%%%%%%%%%%%%%%%%%%%%%%%
% \appendix
% % ==============================================================================
% % APPENDIX: STRUCTURE AND TODO LIST
% % ==============================================================================
% % TODO LIST FOR APPENDICES:
% % - [  ] Write complete geometric theory (ray-fibers, immersion, injectivity)
% % - [  ] Write implementation details (network architectures, training, hyperparameters)
% % - [  ] Add additional experiments (ablations, sensitivity, failure cases)
% % - [  ] Write EVT background (regular variation, tail measures, Hill estimator)
% % ==============================================================================

% \section{Complete Geometric Theory}
% \label{app:geometry}

% % TODO: Write complete ray-fiber theory
% \subsection{Ray-Fiber Analysis}

% TODO: Develop complete geometric theory. Include:
% \begin{itemize}
%     \item Ray-fibers \(\mathcal{F}_c(u) = \cM \cap R_{c,u}\) where \(R_{c,u} = \{c + tu : t > 0\}\)
%     \item Radial-value fibers \(\mathcal{S}_c(u)\)
%     \item Quotient sets \(\mathcal{Q}_c(u,v) = \mathcal{S}_c(u) / \mathcal{S}_c(v)\)
%     \item Hierarchy of fiber complexity (singleton, finite, discrete, positive-dimensional)
% \end{itemize}

% % TODO: Write immersion criteria
% \subsection{Immersion Criteria}

% TODO: State and prove conditions for encoder to be an immersion. Include:
% \begin{itemize}
%     \item Differential analysis: when is \(d(f_c)_x\) injective?
%     \item Angular kernel space \(K_x = d\pi_c(T_x \cM) \cap \text{Ker}(de_c)\)
%     \item Sufficient conditions involving \(e_c\) and \(a_c\)
%     \item Connection to three-component decoder: correction provides needed degrees of freedom
% \end{itemize}

% % TODO: Write injectivity conditions
% \subsection{Injectivity and Global Collisions}

% TODO: Analyze when encoder is globally injective. Include:
% \begin{itemize}
%     \item Collision criterion: \(f_c(x_1) = f_c(x_2) \iff e_c(u_1) = e_c(u_2)\) and \(a_c(u_2)/a_c(u_1) \in \mathcal{Q}_c(u_1, u_2)\)
%     \item Equal-\(e_c\) fibers: \(\mathcal{E}_\xi = \{u : e_c(u) = \xi\}\)
%     \item Singleton-fiber simplification
%     \item Role of magnitude component \(b\) in avoiding collisions
% \end{itemize}

% \section{Implementation Details}
% \label{app:implementation}

% % TODO: Write network architectures
% \subsection{Network Architectures}

% TODO: Specify all architectural details. Include:
% \begin{itemize}
%     \item Encoder: \texttt{encoder\_angular} (3 → 128 → 128 → 2, ReLU, then normalize), \texttt{encoder\_magnitude} (3 → 128 → 128 → 1, ReLU + softplus)
%     \item Decoder: \texttt{decoder\_angular\_base} (2 → 128 → 128 → 3, ReLU, then normalize), \texttt{decoder\_angular\_correction} (3 → 128 → 128 → 3, ReLU), \texttt{decoder\_magnitude} (2 → 128 → 128 → 1, ReLU + softplus)
%     \item Initialization: Xavier/Glorot for weights, zero for biases
%     \item Total parameters: [TODO: count]
% \end{itemize}

% % TODO: Write training procedure
% \subsection{Training Procedure}

% TODO: Document training details. Include:
% \begin{itemize}
%     \item Optimizer: Adam with \(\beta_1 = 0.9\), \(\beta_2 = 0.999\)
%     \item Learning rate: \(10^{-3}\) with cosine annealing schedule
%     \item Batch size: 256 (synthetic), 128 (real-world)
%     \item Epochs: 1000 (synthetic), 500 (real-world)
%     \item Loss weights: \(\lambda_{\text{ang}} = 1.0\), \(\lambda_{\text{rad}} = 1.0\), \(\lambda_{\text{sparse}} \in \{0.001, 0.01, 0.1\}\)
% \end{itemize}

% % TODO: Write hyperparameter ranges
% \subsection{Hyperparameter Selection}

% TODO: Document hyperparameter search. Include:
% \begin{itemize}
%     \item \(p\) (homogeneity degree): tested \(\{1.0, 1.5, 2.0\}\), selected \(p = 2.0\)
%     \item \(\lambda_{\text{sparse}}\): grid search over \(\{0.001, 0.01, 0.1\}\)
%     \item Network depth/width: ablated \(\{64, 128, 256\}\) hidden units
%     \item Center \(c\): fixed at data centroid (could be learned)
% \end{itemize}

% % TODO: Write computational requirements
% \subsection{Computational Requirements}

% TODO: Document compute resources. Include:
% \begin{itemize}
%     \item Hardware: [TODO: specify GPU type, RAM]
%     \item Training time: [TODO: hours per experiment]
%     \item Inference time: [TODO: ms per sample]
%     \item Code: PyTorch 2.x, available at [TODO: github link]
% \end{itemize}

% \section{Additional Experiments}
% \label{app:experiments}

% % TODO: Write additional ablations
% \subsection{Additional Ablations}

% TODO: Report additional ablation studies. Include:
% \begin{itemize}
%     \item Effect of \(p\) on tail preservation
%     \item Effect of network depth/width
%     \item Effect of center choice \(c\)
%     \item Sensitivity to loss weight ratios
% \end{itemize}

% % TODO: Write sensitivity analyses
% \subsection{Sensitivity Analysis}

% TODO: Analyze robustness. Include:
% \begin{itemize}
%     \item Robustness to initialization
%     \item Robustness to batch size
%     \item Robustness to noise in data
%     \item Cross-dataset generalization
% \end{itemize}

% % TODO: Write failure cases
% \subsection{Failure Cases}

% TODO: Document when method fails. Include:
% \begin{itemize}
%     \item Manifolds with persistent curvature (no cone-like tail)
%     \item Very high-dimensional ambient space (\(D \gg 3\))
%     \item Non-regularly-varying data
%     \item Visualizations of failure modes
% \end{itemize}

% \section{Extreme Value Theory Background}
% \label{app:evt}

% % TODO: Write EVT background
% \subsection{Regular Variation}

% TODO: Define regular variation precisely. Include:
% \begin{itemize}
%     \item Definition: \(t P(X/t \in A) \to \mu(A)\) as \(t \to \infty\)
%     \item Examples: Pareto, Student-\(t\), Fr\'echet distributions
%     \item Connection to max-domain of attraction
%     \item Tail index \(\alpha\): \(P(\|X\| > t) \sim t^{-\alpha}\)
% \end{itemize}

% % TODO: Write tail measures
% \subsection{Tail Measures and Pushforwards}

% TODO: Explain tail measures. Include:
% \begin{itemize}
%     \item Tail measure \(\mu\) on \(\RR^D \setminus \{0\}\)
%     \item Scaling property: \(\mu(tA) = t^{-\alpha} \mu(A)\) for \(t > 0\)
%     \item Pushforward \(f_* \mu\): \((f_* \mu)(B) = \mu(f^{-1}(B))\)
%     \item Key result: homogeneous maps preserve regular variation
% \end{itemize}

% % TODO: Write Hill estimator
% \subsection{Hill Estimator}

% TODO: Document Hill estimator for \(\alpha\). Include:
% \begin{itemize}
%     \item Definition: \(\hat{\alpha} = \left( \frac{1}{k} \sum_{i=1}^k \log(X_{(i)} / X_{(k+1)}) \right)^{-1}\)
%     \item Choice of \(k\) (number of order statistics)
%     \item Bias-variance trade-off in \(k\)
%     \item Implementation details and confidence intervals
% \end{itemize}

% ##############################################################################
% ##  CHECKLIST (NeurIPS 2026) -- actioned at migration A0.4                  ##
% ##############################################################################
% TODO CHK-0: Do NOT paste the checklist into the 2025_report tex. It belongs
% in the 2026_report/neurips_2026.tex migration target (file already exists
% at 2026_report/checklist.tex -- \input it in the 2026 document per template).
% TODO CHK-1: When migrating, fill each \answerTODO / \justificationTODO per
% MASTER PLAN §A7. Pre-drafted answers:
%   Q1 Claims       Yes  -- "Abstract and §1 state exactly the two theoretical
%                            results (Prop 1, Thm 2) and the headline
%                            experimental numbers, all of which are proved
%                            / reported in the body."
%   Q2 Limitations  Yes  -- "§7 Limitations (L1 single-chart, L2 Lipschitz
%                            rate, L3 Hill k-selection, L4 m << D regime)."
%   Q3 Theory       Yes  -- "Prop 1 stated in §4.1 with proof in App A; Thm 2
%                            stated in §4.2 with proof in App A."
%   Q4 Repro        Yes  -- "App D specifies architectures, schedule, metrics;
%                            code at <anonymised URL> reproduces Fig 3, 4, 5."
%   Q5 Data/code    Yes  -- "Anonymised Zenodo at submission; GitHub at
%                            camera-ready. ERA5 from Copernicus CDS."
%   Q6 Exp settings Yes  -- "§5 splits + §D full hyperparameters."
%   Q7 Stat sig     Yes  -- "Table 1 reports mean +/- std over 3+ seeds
%                            (TODO: re-run exp02 with num_seeds >= 3)."
%   Q8 Compute      Yes  -- "App D Compute subsection: hardware, wall-clock,
%                            total GPU-hours."
%   Q9 Ethics       Yes  -- "Conforms; no human subjects; EVT on public
%                            reanalysis data."
%   Q10 Broad imp   Yes  -- "§7 Broader Impact; positive applications in
%                            climate risk / finance risk."
%   Q11 Safeguards  NA   -- "Dimension-reduction method; no generative model
%                            of dual-use concern."
%   Q12 Licenses    Yes  -- "ERA5 Copernicus Open License (Hersbach_2020);
%                            PyTorch BSD; matplotlib PSF-like; all cited."
%   Q13 New assets  Yes  -- "Released code + trained checkpoints, documented
%                            in App D."
%   Q14 Human subj  NA   -- "No human-subjects research."
%   Q15 IRB         NA   -- "No human-subjects research."
%   Q16 LLM usage   NA   -- "LLMs not used in the methodology. (If assistive
%                            writing was used, state so here.)"
% ##############################################################################

% %%%%%%%%%%%%%%%%%%%%%%%%%%%%%%%%%%%%%%%%%%%%%%%%%%%%%%%%%%%%
% \newpage
% \section*{NeurIPS Paper Checklist}

% \begin{enumerate}

% \item {\bf Claims}
%     \item[] Question: Do the main claims made in the abstract and introduction accurately reflect the paper's contributions and scope?
%     \item[] Answer: \answerTODO{}
%     \item[] Justification: \justificationTODO{}

% \item {\bf Limitations}
%     \item[] Question: Does the paper discuss the limitations of the work performed by the authors?
%     \item[] Answer: \answerTODO{}
%     \item[] Justification: \justificationTODO{}

% \item {\bf Theory assumptions and proofs}
%     \item[] Question: For each theoretical result, does the paper provide the full set of assumptions and a complete (and correct) proof?
%     \item[] Answer: \answerTODO{}
%     \item[] Justification: \justificationTODO{}

% \item {\bf Experimental result reproducibility}
%     \item[] Question: Does the paper fully disclose all the information needed to reproduce the main experimental results of the paper to the extent that it affects the main claims and/or conclusions of the paper (regardless of whether the code and data are provided or not)?
%     \item[] Answer: \answerTODO{}
%     \item[] Justification: \justificationTODO{}

% \item {\bf Open access to data and code}
%     \item[] Question: Does the paper provide open access to the data and code, with sufficient instructions to faithfully reproduce the main experimental results, as described in supplemental material?
%     \item[] Answer: \answerTODO{}
%     \item[] Justification: \justificationTODO{}

% \item {\bf Experimental setting/details}
%     \item[] Question: Does the paper specify all the training and test details (e.g., data splits, hyperparameters, how they were chosen, type of optimizer, etc.) necessary to understand the results?
%     \item[] Answer: \answerTODO{}
%     \item[] Justification: \justificationTODO{}

% \item {\bf Experiment statistical significance}
%     \item[] Question: Does the paper report error bars suitably and correctly defined or other appropriate information about the statistical significance of the experiments?
%     \item[] Answer: \answerTODO{}
%     \item[] Justification: \justificationTODO{}

% \item {\bf Experiments compute resources}
%     \item[] Question: For each experiment, does the paper provide sufficient information on the computer resources (type of compute workers, memory, time of execution) needed to reproduce the experiments?
%     \item[] Answer: \answerTODO{}
%     \item[] Justification: \justificationTODO{}
    
% \item {\bf Code of ethics}
%     \item[] Question: Does the research conducted in the paper conform, in every respect, with the NeurIPS Code of Ethics \url{https://neurips.cc/public/EthicsGuidelines}?
%     \item[] Answer: \answerTODO{}
%     \item[] Justification: \justificationTODO{}

% \item {\bf Broader impacts}
%     \item[] Question: Does the paper discuss both potential positive societal impacts and negative societal impacts of the work performed?
%     \item[] Answer: \answerTODO{}
%     \item[] Justification: \justificationTODO{}
    
% \item {\bf Safeguards}
%     \item[] Question: Does the paper describe safeguards that have been put in place for responsible release of data or models that have a high risk for misuse (e.g., pretrained language models, image generators, or scraped datasets)?
%     \item[] Answer: \answerTODO{}
%     \item[] Justification: \justificationTODO{}

% \item {\bf Licenses for existing assets}
%     \item[] Question: Are the creators or original owners of assets (e.g., code, data, models), used in the paper, properly credited and are the license and terms of use explicitly mentioned and properly respected?
%     \item[] Answer: \answerTODO{}
%     \item[] Justification: \justificationTODO{}

% \item {\bf New assets}
%     \item[] Question: Are new assets introduced in the paper well documented and is the documentation provided alongside the assets?
%     \item[] Answer: \answerTODO{}
%     \item[] Justification: \justificationTODO{}

% \item {\bf Crowdsourcing and research with human subjects}
%     \item[] Question: For crowdsourcing experiments and research with human subjects, does the paper include the full text of instructions given to participants and screenshots, if applicable, as well as details about compensation (if any)?
%     \item[] Answer: \answerNA{}
%     \item[] Justification: The current paper is theoretical and geometric; no crowdsourcing or human-subject research is presently involved.

% \item {\bf Institutional review board (IRB) approvals or equivalent for research with human subjects}
%     \item[] Question: Does the paper describe potential risks incurred by study participants, whether such risks were disclosed to the subjects, and whether Institutional Review Board (IRB) approvals (or an equivalent approval/review based on the requirements of your country or institution) were obtained?
%     \item[] Answer: \answerNA{}
%     \item[] Justification: The current paper does not involve human-subject research.

% \item {\bf Declaration of LLM usage}
%     \item[] Question: Does the paper describe the usage of LLMs if it is an important, original, or non-standard component of the core methods in this research?
%     \item[] Answer: \answerNA{}
%     \item[] Justification: The current research problem and methodology do not use LLMs as a component of the proposed theory or method.

% \end{enumerate}


